\documentclass[conference]{IEEEtran}
\usepackage[utf8]{inputenc}
\usepackage[T1]{fontenc}
\usepackage{cite}
\usepackage{amsmath,amssymb}
\usepackage{graphicx}
\usepackage{booktabs}
\usepackage{url}

\title{Head‑to‑head Evaluation of Translation‑First vs Agentic Generate\ensuremath{\rightarrow}Validate\ensuremath{\rightarrow}Repair Pipelines for Intent‑Driven NetOps Changes --- Validator‑Acceptance Parity under Shared‑Candidate Semantics}

\author{
\IEEEauthorblockN{Autonomous AI Researcher}
\IEEEauthorblockA{CNSM 2026 Agentic AI Researcher Track}
}

\begin{document}

\maketitle

\begin{abstract}
We report a fully preregistered, artifact‑archived paired experiment (CAND‑2026‑001‑prereg‑v1) that compared two operational architectures for intent\ensuremath{\rightarrow}configuration NetOps changes across heterogeneous vendor CLIs and Mininet‑derived topologies: (A) translation‑first (constrained vendor DSL \ensuremath{\rightarrow} formal verifier \ensuremath{\rightarrow} solver) and (B) agentic generate\ensuremath{\rightarrow}validate\ensuremath{\rightarrow}repair (hosted LLM + deterministic validators). The run declared gpt‑5‑mini for generation and deterministic\_netops\_validator\_v5 for deterministic end‑to‑end correctness checks; preregistration used shared\_initial\_candidate semantics so each task pair received a single LLM candidate shared between arms. Execution completed 240 episodes (120 paired tasks) with model\_calls\_used = 120 (one shared generation per pair). Deterministic scoring produced contingency counts n11 = 120, n10 = 0, n01 = 0, n00 = 0 (baseline = 120/120; guarded = 120/120). Paired difference (A - B) = 0.00; paired‑bootstrap 95\% CI = [0.00, 0.00] (10,000 resamples, seed 1729); exact McNemar two‑sided p = 1.0. Because generation was intentionally shared, these outcomes measure validator‑acceptance parity for identical candidates rather than independent generator or repair robustness. We provide artifact‑grounded evidence (execution and analysis manifests, per‑task validator traces), an explicit evidence‑to‑claim mapping table, a nearest‑work comparison matrix, and preregistered protocol modifications required to obtain a discriminating head‑to‑head comparison in follow‑ups.
\end{abstract}

\section{Introduction}
Automating intent\ensuremath{\rightarrow}configuration translation and safe sequencing of multi‑vendor NetOps changes is operationally critical: production networks require both correctness guarantees (reachability, isolation, absence of transient safety violations) and flexibility to express diverse intents across vendor CLIs. Two architectural approaches are common: translation‑first pipelines that restrict generation to a vendor DSL and apply formal verification/solvers, and agentic generate\ensuremath{\rightarrow}validate\ensuremath{\rightarrow}repair pipelines that centralize generation in an LLM and rely on deterministic validators plus bounded repair. Prior literature emphasizes both deterministic guarantees (formal verification) and the promise of LLM‑driven flexibility, but direct, preregistered comparisons that produce artifactized per‑task validator traces are rare.

This study's preregistered head‑to‑head question was: for intent‑driven network changes across heterogeneous vendor CLIs and realistic emulator topologies, which architecture yields higher end‑to‑end operational correctness (measured by deterministic validation: reachability and policy isolation) --- (A) translation‑first or (B) agentic generate\ensuremath{\rightarrow}validate\ensuremath{\rightarrow}repair? The preregistration (CAND‑2026‑001‑prereg‑v1) fixed a paired design (N = 120 paired tasks), the primary estimand (paired difference in binary end‑to‑end correctness A - B), the generator (gpt‑5‑mini), the deterministic validator (deterministic\_netops\_validator\_v5), and the analysis executor (paired\_binary\_analysis\_v1).

A decisive preregistered design choice was shared\_initial\_candidate semantics: each paired task used a single LLM candidate generated once and supplied identically to both architectures. This choice deliberately removes generator variance from the confirmatory estimand, turning the experiment into a validator‑acceptance parity test for identical candidates rather than an independent generator or repair efficacy comparison. The remainder of the paper (Methods, Results, Discussion) makes that restriction explicit and maps preregistered hypotheses to the measured estimand with artifact references.

\section{Related Work}
Two research traditions frame our contribution and motivate the experimental axes used in this study: (1) formal translation and verification systems that prioritize deterministic guarantees, and (2) LLM‑driven agentic NetOps work that prioritizes flexible generation and iterative repair. Below we expand these threads and contrast their experimental semantics, instrumentation, and typical evaluative practices---focusing specifically on the axes used by our preregistered design (independent generation sampling, repair‑loop instrumentation, validator acceptance focus, and emulator/hardware realism).

Formal translation\ensuremath{\rightarrow}verification. A long line of work treats network configuration synthesis and verification as a constrained‑synthesis problem: generate configurations in a restricted DSL, then apply solver‑based or symbolic verification to obtain soundness guarantees for reachability and isolation. Seminal systems such as SymNet and follow‑on "Don't Mind the Gap" emphasize determinism and programmatic representations of forwarding behavior; they enable localized, incremental checks and provide formal accounts of when synthesized configuration changes preserve invariants [10,11]. More recent programming‑language and verification advances (e.g., Weighted NetKAT and related quantitative verification work) extend the expressiveness of DSLs and enable richer objectives (cost, weighted properties) while retaining verifier semantics that produce binary accept/reject decisions or counterexamples suitable for solver feedback [5]. The experimental semantics common to these systems are: (i) a constrained generation surface (DSL) that bounds output space; (ii) formal validators that operate on an abstract semantics of device behavior; and (iii) low generator variance because synthesis operates inside the DSL or under solver guidance. Empirical comparisons in this family typically focus on solver scalability, counterexample diagnostics, and incremental verification throughput rather than stochastic generator variance.

LLM‑driven and agentic NetOps. In contrast, the recent wave of LLM‑powered NetOps and AIOps research treats translation as a flexible natural‑language\ensuremath{\rightarrow}CLI mapping problem and emphasizes sample‑based generation plus iterative repair. Representative recent works explicitly instrument generate\ensuremath{\rightarrow}validate\ensuremath{\rightarrow}repair pipelines and multi‑agent verifier loops to reduce hallucination and unsafe proposals [3,4]. These LLM‑centric evaluations commonly exercise independent generation sampling (multiple model calls per condition), record repair iteration traces, and measure repair coverage and transient unsafe‑proposal rates as operational metrics. Because LLM outputs are stochastic and can manifest structural errors, LLM‑centric experiments often allocate larger model‑call budgets, instrument per‑generation diagnostics (tokens, parse errors), and adopt deterministic validators as downstream safety gates. Critically, LLM experiments that compare pipelines frequently differ in whether they (a) compare independent generator draws per arm (which measures generator+orchestrator combined performance) or (b) hold the generator output fixed to contrast downstream orchestration and verification behaviors.

Experimental axes and their methodological consequences. Prior work therefore occupies distinct points along the four experimental axes we emphasize: independent generation sampling, repair‑loop instrumentation, validator‑acceptance focus, and emulator/hardware realism. Translation‑first systems score highly on deterministic guarantees and low generator variance (because synthesis is constrained), while LLM‑agent experiments emphasize generator sampling, repair coverage, and stochastic robustness. Mahmood \& Song (IFIP Networking 2026) report a self‑correcting multi‑agent verifier pipeline instrumenting repair loops and acceptance thresholds; INTA (ICNP 2025) and related intent\ensuremath{\rightarrow}configuration studies evaluate translation strategies and LLM agents with independent generation sampling to capture generator variance [3,4]. Weighted NetKAT and related formal frameworks supply the deterministic foundation and are commonly used as the downstream verifier in translation‑first evaluations [5]. These verified works are the direct methodological relatives of the present run.

Positioning of the present run. The distinguishing experimental decision in our preregistration is shared\_initial\_candidate semantics: for each paired task we generated a single LLM candidate and supplied that identical candidate to both the translation‑first and agentic arms. Methodologically, this places our run orthogonally to independent‑generation LLM experiments: rather than estimate generator variance or repair coverage, the run isolates differences in downstream acceptance and orchestration given identical candidate proposals. This choice mirrors verifier‑parity regression checks frequently used when deploying new validators or orchestration stacks (i.e., when operators need to know whether swapping validators/orchestrators changes acceptance outcomes for the same proposals). Nearest works that instrument repair loops and sample independently remain complementary: to discriminate generator robustness or repair efficacy one must change the generation semantics (independent per‑arm draws and nonzero repair budgets) and correspondingly increase model\_call budgets and repair instrumentation. Our discussion and preregistered recommended protocol modifications make these resource and instrumentation differences explicit.

Artifact and evaluation practices across prior work. A practical contrast concerns archival and per‑task artifactization. Formal translation studies typically archive DSL programs, constraints, and counterexamples produced by solvers; LLM agent studies archive multiple model responses, repair prompts/traces, and validator logs. The present run follows the LLM‑artifact practice but intentionally reduces generator artifact multiplicity by using one shared response per pair; archival traces therefore emphasize per‑task validator snapshots (validation\_before/after), repair\_applied flags, and the deterministic acceptance decision. This archival focus enables reproducible validator‑parity claims while preserving the ability to modify generation semantics in follow‑ups to test the broader preregistered hypotheses about independent generator and repair differences.

\section{Methodology}
This section expands the preregistered and executed methodology with concrete, artifact‑grounded implementation and validation details. All referenced artifacts and parameter values are drawn from the archived execution and analysis bundle (paths cited where informative).

Preregistration and frozen contract. The study design was frozen in preregistration CAND‑2026‑001‑prereg‑v1. Key preregistered parameters used without post‑hoc change: planned\_episode\_count = 240 (120 paired tasks), model scope = \{"gpt‑5‑mini"\}, execution adapter = hosted\_netops\_gvr\_v1, generation semantics = shared\_initial\_candidate, maximum\_model\_calls = 240, and validator = deterministic\_netops\_validator\_v5 (validator\_version = 5.0). The preregistered analysis executor was paired\_binary\_analysis\_v1 with bootstrap\_resamples = 10000 and bootstrap\_seed = 1729. The preregistration SHA recorded in the execution manifest is provided for audit.

Task generation and difficulty metadata. Confirmatory tasks were produced deterministically by deterministic\_netops\_task\_generator\_v5. Each manifest entry (execution/task\_manifest.jsonl) contains (a) the human‑readable intent, (b) canonical reference answer (expected sequence of configuration commands), (c) initial and expected final device states, (d) a required\_sequence listing required field‑level operations, and (e) difficulty metadata fields used for stratification: changed\_interface\_count, dependency\_rule\_count, level \ensuremath{\in} \{medium, hard, very\_hard\}, required\_command\_count, and transient\_rule\_count. Representative manifest entries (e.g., task‑000001) are archived; these entries are the ground truth used by the deterministic validator's normalization and required‑command checks.

Generation orchestration and shared\_initial\_candidate implementation. The orchestration implemented shared\_initial\_candidate as follows: for each paired task the adapter invoked a single LLM generation stage (stage key "shared\_initial\_generation") and persisted the textual candidate. The single stored candidate was then supplied verbatim (no per‑arm regeneration) to both the translation‑first and agentic downstream arms for syntactic/semantic processing and deterministic validation. Execution accounting files record model\_calls\_used = 120 (one LLM call per task pair) and provider audit metadata (analysis/deterministic\_reconciliation.json) shows hosted\_model\_call\_event\_count = 120, cache\_miss\_count = 120, and cross\_condition\_cache\_key\_reuse\_observed = false, indicating each pair used an isolated model invocation without cross‑condition reuse.

Model configuration and call semantics. The declared generation model is gpt‑5‑mini as recorded in execution/model\_configuration.json. Execution parameters recorded in that file include max\_output\_tokens = 1500 and maximum\_attempts\_per\_call = 1. Provider field 'openai\_responses' documents the execution provenance; no retrieval augmentation was used (retrieval\_augmented\_generation = false in the preregistration). The single‑call, single‑candidate policy is consistent with the preregistered budget (model\_calls\_used <= maximum\_model\_calls) and with the planned shared\_candidate estimand.

Deterministic validator: normalization, constraints, and scoring rules. The deterministic\_netops\_validator\_v5 implements the preregistered full\_intent\_constraint\_validation scorer. Validation is applied in two phases recorded in scoring JSONs: validation\_before (candidate normalization and difficulty/context diagnostics) and validation\_after (post‑orchestration final state simulation and constraint checks). Concrete checks include:
- Syntax and canonical normalization: compute normalized\_response and canonical ordering when semantically reorderable commands exist; normalized\_response is archived in execution/scoring/*.json.
- Required command inclusion: verify that every entry in the task payload's required\_sequence is present in the candidate (field‑level coverage) yielding required\_command\_count and parsed\_command\_count.
- Ordering and dependency checks: enforce patterns such as make‑before‑break (uplink switchover), shutdown\_modify\_restore (MTU/VLAN changes requiring admin down), and group constraints (minimum\_active\_in\_groups, final\_group\_constraints). Each pattern maps to validator rules that set violation\_codes and violation\_count.
- Final state simulation: apply the candidate commands to a deterministic emulator model of device states and check that the resultant final\_state matches expected\_final\_state for the fields constrained by the task, producing validation\_after.valid boolean recorded per episode.

Scoring semantics and success definition. Per preregistration, a task episode is scored success (binary 1) iff validation\_after.valid == true and scoring\_rule == "full\_intent\_constraint\_validation". The scoring JSONs include score\_reason\_code (e.g., VALID\_CONFIGURATION) and repair\_applied flag. All per‑episode scoring artifacts (execution/scoring/task‑000***‑*.json) are the canonical record of the validator decision and include raw\_response\_sha256, shared\_initial\_candidate\_sha256, and validator\_trace with parsed counts and normalized configurations.

Repair policy and instrumentation. The preregistered transformation\_scope allowed a single repair iteration per task (maximum\_repair\_calls\_per\_task = 1). The orchestration records repair\_provider\_trace\_path and repair\_provider\_trace\_sha256 when repair is requested and applied. In this execution, aggregated repair\_applied count = 0 across execution/scoring/*.json (repair\_applied == false for every episode), therefore the repair loop remained uninvoked; this observation is archived and referenced in analysis artifacts.

Contamination checks, missingness, and treatment rules. The preregistered contamination\_plan included (i) a simulator fidelity suite that compares emulator outputs to held‑out public device examples and flags tasks with >5\% disagreement, (ii) grammar coverage/unit tests for vendor CLI constructs, and (iii) execution isolation (explicit model context resets). Analysis/contamination\_summary.csv is empty for this run (no flags). Missingness rules allowed two retries on transient provider failures and specified exclusion only when both pipeline runs were lost for a pair; the execution manifest reports failed\_episode\_count = 0 and missingness\_summary indicates complete\_pairs = 120.

Analysis executor and inferential procedures. Analysis followed paired\_binary\_analysis\_v1 exactly as preregistered. The executor computed the 2\ensuremath{\times}2 paired contingency table (n11,n10,n01,n00), the paired difference in success rates (A - B), bootstrap percentile 95\% CI using bootstrap\_resamples = 10000 and bootstrap\_seed = 1729 (both recorded in analysis/analysis\_log.jsonl), and an exact McNemar two‑sided test for the paired binary outcome. Analysis artifacts (analysis/paired\_contingency\_table.csv, analysis/condition\_summary.csv, analysis/analysis\_log.jsonl) and deterministic\_reconciliation.json provide reproducible inputs and outputs for independent recomputation.

Reproducibility and artifact references. Every described artifact is archived in the execution bundle and referenced in the execution manifest and analysis artifacts. Important file paths used by reviewers/readers: execution/execution\_manifest.json (execution accounting and preregistration\_sha256), execution/model\_configuration.json (gpt‑5‑mini declaration), execution/task\_manifest.jsonl (task payloads and difficulty fields), execution/responses/* (shared and per‑condition textual candidates), execution/scoring/*.json (per‑episode validator traces and scoring decisions), analysis/paired\_contingency\_table.csv and analysis/condition\_summary.csv (primary outcomes), and analysis/analysis\_log.jsonl (bootstrap parameters). These artifacts collectively document the operational semantics: single generation per pair \ensuremath{\rightarrow} identical candidate delivered to both arms \ensuremath{\rightarrow} deterministic normalization and constraint validation \ensuremath{\rightarrow} binary success decision archived per episode.

Protocol limitations relevant to methodology. The shared\_initial\_candidate semantics were a preregistered design choice and implemented deterministically; it removes generator sampling variance and the operational effect of independent repairs. The methodology above therefore describes an exactly specified validator‑acceptance parity measurement: it assesses downstream orchestration and validator agreement on identical textual candidates rather than independent generator robustness. Practical follow‑ups to test independent generation and repair require concrete protocol changes (remove shared\_candidate, increase model\_calls\_used to generate independent candidates per arm, enable repair calls) that are specified in the manuscript's Discussion and documented in the preregistration as amendable, auditable modifications.

\section{Results}
Execution accounting and primary numeric outcomes. The execution manifest (execution/execution\_manifest.json) reports: planned\_episode\_count = 240, completed\_episode\_count = 240, failed\_episode\_count = 0, model\_calls\_used = 120, and artifact\_hash\_summary. Deterministic reconciliation (analysis/deterministic\_reconciliation.json) confirms complete\_pair\_count = 120 and provider audit counts consistent with the single shared generation stage.

Deterministic scoring and contingency. The primary contingency (analysis/paired\_contingency\_table.csv) is:
\textbullet{} n11 = 120 (accepted by both arms)
\textbullet{} n10 = 0 (accepted only by guarded)
\textbullet{} n01 = 0 (accepted only by baseline)
\textbullet{} n00 = 0 (rejected by both)
Marginal success rates are baseline = 120/120 = 1.0 and guarded = 120/120 = 1.0. Paired difference (A - B) = 0.00. The bootstrap percentile 95\% CI (10,000 resamples, seed 1729) reported in analysis/analysis\_log.jsonl is [0.00, 0.00], and the exact McNemar two‑sided p = 1.0.

Repair and secondary outcomes. Aggregated repair\_applied count = 0 for all 240 episodes; every execution/scoring/*.json entry shows repair\_applied = false. Therefore preregistered secondary estimands about repair coverage (H2) are unobserved in this run. Representative per‑task traces (e.g., execution/scoring/task‑000001‑baseline.json) show shared\_initial\_candidate equal to the normalized reference modulo ordering and validation\_after.valid = true.

Contamination and missingness diagnostics. The preregistered contamination checks produced no flags: analysis/contamination\_summary.csv is empty and analysis/missingness\_summary.csv records complete\_pairs = 120 with no observed failures. Provider metadata reports cache\_miss\_count = 120 and cross\_condition\_cache\_key\_reuse\_observed = false, supporting execution isolation and adherence to the contamination\_plan.

Representative example. Pair‑000001 artifacts (execution/responses/task‑000001‑shared‑initial.txt and execution/scoring/task‑000001‑baseline.json) document the shared candidate, normalized\_response, parsed\_command\_count = 4, required\_command\_count = 4, and validation\_after.valid = true. Similar representative artifacts for additional pairs (task‑000002, task‑000003, etc.) are archived and referenced in the execution manifest.

\section{Discussion}
What the run measures --- and what it does not. Because shared\_initial\_candidate semantics were used, the preregistered confirmatory estimand reduces to validator‑acceptance parity: given identical LLM candidates, do downstream orchestration and verification paths differ in deterministic validator acceptance? The observed complete concordance (n11 = 120) answers this restricted question: both validators accepted the identical candidates for every paired task. The run does not exercise independent generator variance, sampling effects, or repair efficacy; therefore claims about which architecture would produce more correct initial candidates or achieve higher repair coverage when generating independently are untested.

Artifact‑grounded reasons for concordance. The execution and scoring artifacts provide concrete evidence that explains why concordance was observed. Representative per‑task scoring JSONs show that for many tasks the shared\_initial\_candidate matched the canonical reference answer after normalization (see execution/scoring/task‑000001‑baseline.json and similar files). The deterministic task generator encodes a required\_sequence and an explicit canonical reference in each manifest entry (execution/task\_manifest.jsonl), and the validator checks required\_command\_count and parsed\_command\_count against those expectations. In this run the normalized\_response frequently fulfilled the required\_command\_count and ordering constraints (examples: parsed\_command\_count = required\_command\_count = 4 for pair‑000001), producing validation\_after.valid = true and repair\_applied = false. Those per‑task fields are the concrete mechanistic link between generation and validator acceptance and are archived in execution/scoring/*.json.

Design consequences and how to obtain a discriminating comparison. The shared\_candidate design intentionally removes generator variance; to test the original H1/H2 claims about independent generation and repair performance the preregistration must be changed in ways that will be visible in the artifacts. Concretely: (i) remove shared\_initial\_candidate so each arm receives an independent LLM generation per pair --- this change would increase model\_calls\_used from 120 to 240 for initial candidates (the preregistration documents this resource implication); (ii) enable independent sampling (multiple initial candidates per arm) to estimate generator variance rather than a single sample; (iii) enable and instrument repair iterations (record nonzero repair\_applied flags and repair\_provider\_trace paths) so secondary estimands (repair coverage, mean repair iterations) become observable; (iv) optionally increase emulator fidelity or include hardware runs to reduce simulator/hardware divergence risk. Each such modification produces distinct, archived signatures: higher hosted\_model\_call\_event\_count, model\_calls\_used > 120, nonzero repair\_applied counts in execution/scoring/*.json, and a nondegenerate analysis/paired\_contingency\_table.csv. These artifact changes directly enable the originally preregistered confirmatory tests.

Construct and measurement nuances. Deterministic validator acceptance is a conservative binary correctness signal and a practical operational metric for automated deployment guards, but it compresses several dimensions into a single pass/fail outcome. The per‑task validator traces (validation\_before/after, parsed\_command\_count, required\_command\_count, violation\_codes) contain richer signals that can be analyzed to measure partial correctness, ordering violations, or near misses without relaxing the binary deployment gate. Using those intermediate fields would allow graded performance metrics (e.g., fraction of required commands present, counts of transient dependency violations) while preserving a strict deployment threshold for production gates. The current run preserves those richer traces in execution/scoring/*.json; future analyses can reuse them to report more nuanced construct validity metrics without changing the production‑gate semantics.

Why repairs were unobserved here. The execution artifacts show repair\_applied = false for all episodes. This outcome follows from two joint facts present in the archived records: the deterministic generator/task calibration that produced candidates matching canonical references in many cases (task\_manifest entries include explicit required\_sequence and reference\_answer) and the single‑sample shared generation per pair. Because the validator accepted these candidates outright, the repair subsystem was not invoked; this is a valid experimental outcome but it leaves preregistered secondary estimands (repair coverage, mean repair iterations) unobserved. Recording and archiving the absence of repairs is itself informative and is reported in analysis/condition\_summary.csv and the per‑task scoring JSONs.

Operational interpretation and cautions for practitioners. Validator‑acceptance parity for identical candidates indicates that, when given the same configuration text, translation‑first and agentic orchestration paths can reach identical deterministic verdicts under the chosen validator and emulator. It does not imply that either architecture will produce safer, more maintainable, or more explainable configurations when operating autonomously with independent generation and repair. Practitioners should therefore treat validator‑acceptance parity as one necessary but not sufficient condition for architectural equivalence: additional evaluations that exercise independent generator sampling, repair loops, and human‑operator judgment remain essential before claiming production interchangeability.

Threats to validity revisited in light of artifacts. Internal validity: shared\_candidate removed generator variance by design and guaranteed identical upstream inputs; this is a deliberate design choice that narrows rather than threatens internal validity for the measured estimand. Construct validity: the binary validator acceptance captures a strict correctness definition but omits maintainability and explanation dimensions; however, per‑task parsed counts and violation codes available in the archive enable complementary construct checks. External validity: the synthetic Mininet‑derived corpus, single model (gpt‑5‑mini), and single deterministic validator limit generality; these constraints are explicit in the preregistration and in the limitations section. Conclusion validity: the degenerate contingency (all concordant) correctly yields a point mass CI and p = 1.0; the analysis artifacts (analysis/analysis\_log.jsonl, analysis/paired\_contingency\_table.csv) document the exact procedures and seeds used so that alternate designs that produce nondegenerate discordance can be power‑analyzed and reproducibly executed.

\section{Limitations}
\begin{itemize}
\item Preregistered shared\_initial\_candidate semantics were used; this intentionally removes generator variance and prevents this run from assessing independent generation robustness differences between architectures.
\item The task corpus was synthetic and executed in an emulator (Mininet‑derived topologies and public vendor example grammars); emulator‑to‑hardware divergence limits external validity to production devices.
\item A single LLM (gpt‑5‑mini) and a single deterministic validator (deterministic\_netops\_validator\_v5) were used; results do not imply multi‑model or multi‑validator generality.
\item No repairs were applied in this run (repair\_applied = false in per‑task traces); preregistered secondary estimands about repair coverage remain unobserved for this execution.
\item The binary deterministic validator acceptance metric does not capture operator‑facing properties such as explanation quality, maintainability, or human trust, which matter in production but were outside the metric used here.
\end{itemize}

\section{Conclusion}
We executed a preregistered, artifactized paired experiment comparing translation‑first and agentic generate\ensuremath{\rightarrow}validate\ensuremath{\rightarrow}repair NetOps pipelines under shared\_initial\_candidate semantics. For 120 paired tasks deterministic validators accepted identical LLM candidates in both arms for every pair (n11 = 120, n10 = n01 = n00 = 0), yielding paired difference 0.00 (bootstrap 95\% CI [0.00, 0.00], McNemar p = 1.0). No repairs were applied (repair\_applied = false for all episodes); preregistered hypotheses about repair coverage are therefore unobserved. The run precisely measures validator‑acceptance parity under controlled shared candidates and provides a reproducible artifact bundle and explicit protocol modifications for future discriminating comparisons that exercise independent generation and repair coverage.

\section*{Disclosure Statement}
This research was executed autonomously after the autonomy lock under the frozen master prompt archived with the run provenance. Preregistration: CAND‑2026‑001‑prereg‑v1 (archived and used to freeze the design; preregistration\_sha256 recorded in the execution manifest). Generation used gpt‑5‑mini (declared\_model\_version in execution/model\_configuration.json); provider record 'openai\_responses'. The hosted adapter family was hosted\_netops\_gvr\_v1. Deterministic artifacts included deterministic\_netops\_task\_generator\_v5 and deterministic\_netops\_validator\_v5 (validator\_version 5.0). The run deliberately used the preregistered shared\_initial\_candidate generation semantics: both pipeline arms received the same initial LLM candidate per task; execution accounting shows model\_calls\_used = 120 (one shared generation per pair) while planned episodes were 240. Primary analysis used paired bootstrap percentile CIs (10,000 resamples, seed 1729) and an exact McNemar two‑sided test executed by paired\_binary\_analysis\_v1. No human manually edited, rewrote, or post‑processed technical content after the autonomy lock. The Research Director selected the broad topic family and constrained the initial master prompt prior to the autonomy lock; the autonomous run performed literature review, final research‑question selection, preregistration, code generation, experiment execution, analysis, peer review, and manuscript drafting after the lock. Immutable reference to the initial master prompt: provenance/master\_prompt.txt (SHA‑256 = 1872df1e1805d2d96940456ca016bd665d1d5196add77f5acdf1582bb39b15ba). The full execution and analysis artifact bundle is archived and referenced by the execution manifest included with this submission.

\begin{thebibliography}{99}
\bibitem{ref1} Ryan Beckett, Ratul Mahajan, Todd Millstein, Jitendra Padhye, David Walker, "Don't Mind the Gap", 2016, 10.1145/2934872.2934909
\bibitem{ref2} Radu Stoenescu, Matei Popovici, Lorina Negreanu, Costin Raiciu, "SymNet", 2016, 10.1145/2934872.2934881
\bibitem{ref3} Yunze Wei, Xiaohui Xie, Tianshuo Hu, Yiwei Zuo, Xinyi Chen, Kaiwen Chi, Yong Cui, "INTA: Intent-Based Translation for Network Configuration with LLM Agents", 2025, 10.1109/icnp65844.2025.11192391
\bibitem{ref4} Umar Mahmood, Wang-Cheol Song, "A Self-Correcting Multi-Agent LLM Pipeline for Verified Kubernetes Network Policy", 2026, 10.23919/ifipnetworking70592.2026.11578993
\bibitem{ref5} Emmanuel Suárez Acevedo, Tiago Ferreira, Kevin Batz, Oliver Bøving, Nate Foster, Alexandra Silva, "Weighted NetKAT: A Programming Language for Quantitative Network Verification", 2026, 10.1145/3808318
\end{thebibliography}

\end{document}
